% Vietnamese translation for OI Public Library
% Source-PDF: 图论 GRAPH THEORY/经典网络流教程.pdf
\documentclass[11pt]{article}
\usepackage{oipl-vn}

\title{Giáo trình kinh điển về luồng mạng}
\author{Comzyh Blog}
\date{}

\begin{document}
\maketitle

\origtitle{Classic network-flow tutorial}
\sourcepath{Graph Theory / Classic network-flow tutorial.pdf}

\section*{Lời dẫn}

Bản gốc là một bộ slide nhập môn về luồng mạng. Phần dưới đây giữ nguyên nội
dung toán học và các ví dụ chính, nhưng viết lại thành ghi chú dạng bài giảng.
Những hình trong slide được thay bằng mô tả, danh sách cạnh, bảng hoặc công
thức tương đương; các slide ví dụ ở phần cận dưới/cận trên và chi phí nhỏ
nhất đã được kiểm tra trực quan vì chúng chủ yếu là hình mạng.

\section{Một bài toán luồng mạng}

Do con người tiêu hao tài nguyên tự nhiên, người ta nhận ra rằng khoảng sau năm
2300 Trái Đất sẽ không còn thích hợp để cư trú. Vì vậy một khu định cư mới được
xây trên Mặt Trăng để chuẩn bị di dân khi cần. Bất ngờ thay, vào mùa đông năm
2177, do nguyên nhân chưa biết, môi trường Trái Đất sụp đổ dây chuyền; loài
người phải chuyển lên Mặt Trăng trong thời gian ngắn nhất.

Giữa Trái Đất và Mặt Trăng có $n$ trạm vũ trụ, và có $m$ tàu vận tải công cộng
qua lại giữa các trạm đó. Mỗi trạm vũ trụ chứa được vô hạn người, còn tàu thứ
$i$ chỉ chứa được $H_i$ người. Mỗi tàu dừng tuần hoàn tại một dãy trạm. Chẳng
hạn dãy $(1,3,4)$ nghĩa là tàu sẽ dừng theo chu kỳ
$1,3,4,1,3,4,\ldots$. Một tàu đi từ một trạm trong chu kỳ sang trạm kế tiếp
luôn mất $1$ đơn vị thời gian. Người chỉ có thể lên hoặc xuống tàu vào các thời
điểm tàu dừng tại trạm, tại Trái Đất hoặc tại Mặt Trăng. Ban đầu mọi người đều
ở Trái Đất, và mọi tàu đều ở trạm xuất phát. Hãy thiết kế thuật toán tìm phương
án vận chuyển tất cả mọi người lên Mặt Trăng nhanh nhất.

Dữ liệu vào nằm trong \texttt{input.txt}. Dòng đầu gồm ba số nguyên dương
$n,m,k$, lần lượt là số trạm vũ trụ, số tàu và số người cần vận chuyển từ Trái
Đất. Ràng buộc là
\[
  1\le m\le 13,\qquad 1\le n\le 20,\qquad 1\le k\le 50.
\]
Mỗi dòng mô tả một tàu $p_i$: số đầu tiên là sức chứa $H_{p_i}$; số thứ hai là
số trạm trong một chu kỳ dừng $r$, với $1\le r\le n+2$; tiếp theo là
$r$ số hiệu trạm $(S_{i1},S_{i2},\ldots,S_{ir})$. Trái Đất được ký hiệu bằng
$0$, Mặt Trăng bằng $-1$. Tại thời điểm $0$, tất cả các tàu đang ở trạm đầu
tiên của chu kỳ; sau đó tại các thời điểm nguyên $1,2,3,\ldots$ chúng dừng ở
các trạm kế tiếp. Chỉ tại các thời điểm nguyên $0,1,2,\ldots$ mới được lên
hoặc xuống tàu.

Cần ghi ra \texttt{output.txt} thời gian ít nhất để chuyển toàn bộ người đi an
toàn. Nếu bài toán vô nghiệm thì ghi $0$.

\section{Ví dụ đơn giản: luồng cực đại}

Có thể hình dung mạng như một hệ thống ống nước. Trong slide, mỗi cạnh được
ghi dưới dạng $(f,c)$: số bên trái là lưu lượng hiện tại, số bên phải là dung
lượng. Ví dụ mở đầu có nguồn $s$, đích $t$ và bốn đỉnh trung gian
$a,b,c,d$. Lượng luồng hiện tại đi ra khỏi $s$ là $3+2=5$, nên ví dụ đặt câu
hỏi: luồng cực đại có phải bằng $5$ hay không?

\section{Ký hiệu và định nghĩa}

\begin{itemize}
  \item $V$ là tập tất cả các đỉnh trong đồ thị.
  \item $E$ là tập tất cả các cạnh trong đồ thị.
  \item $G=(V,E)$ là toàn bộ đồ thị.
  \item $s$ là nguồn của mạng, $t$ là đích của mạng.
  \item Mỗi cạnh $(u,v)$ có dung lượng $c(u,v)$, trong đó $c(u,v)\ge 0$.
  \item Nếu $c(u,v)=0$ thì xem như cạnh $(u,v)$ không tồn tại trong mạng.
  \item Nếu mạng ban đầu không có cạnh $(u,v)$ thì đặt $c(u,v)=0$.
  \item Mỗi cạnh $(u,v)$ có một lưu lượng $f(u,v)$.
\end{itemize}

\subsection{Ba tính chất của luồng mạng}

Một luồng hợp lệ phải thỏa ba tính chất.

\begin{enumerate}
  \item \term{Giới hạn dung lượng:}
  \[
    f(u,v)\le c(u,v).
  \]
  \item \term{Phản đối xứng:}
  \[
    f(u,v)=-f(v,u).
  \]
  \item \term{Bảo toàn luồng:} với mọi đỉnh không phải nguồn và cũng không
  phải đích, tổng lưu lượng đi vào bằng tổng lưu lượng đi ra.
\end{enumerate}

Kết hợp với phản đối xứng, điều kiện bảo toàn luồng có thể viết là
\[
  \sum_{u\in V} f(v,u)=0
\]
với mọi $v\ne s,t$. Chỉ cần thỏa ba tính chất trên thì ta có một luồng mạng
hợp lệ, còn gọi là \term{luồng khả thi}. Luồng khả thi luôn tồn tại ít nhất là
luồng không.

\subsection{Bài toán luồng cực đại}

Giá trị của một luồng, ký hiệu $|f|$, được định nghĩa là tổng lưu lượng đi ra
từ nguồn:
\[
  |f|=\sum_{v\in V} f(s,v).
\]
Bài toán luồng cực đại yêu cầu cực đại hóa $|f|$ trong khi vẫn thỏa các tính
chất của luồng mạng.

\section{Phân loại cung và mạng dư}

Cho một luồng khả thi $F=(F_{ij})$. Nếu $F_{ij}=C_{ij}$ thì cung đó là
\term{cung bão hòa}; nếu $F_{ij}<C_{ij}$ thì là \term{cung chưa bão hòa}; nếu
$F_{ij}=0$ thì là \term{cung luồng không}; nếu $F_{ij}>0$ thì là
\term{cung có luồng khác không}.

Nếu $P$ là một đường nối nguồn $s$ với đích $t$ trong mạng, không xét hướng của
cạnh, ta quy ước hướng của đường là từ $s$ đến $t$. Khi đó các cung trên
đường chia thành hai loại: cung cùng hướng với đường gọi là \term{cung tiến},
cung ngược hướng với đường gọi là \term{cung lùi}.

\subsection{Mạng dư}

Để cài đặt thuận tiện, ta thường định nghĩa \term{mạng dư} từ mạng ban đầu.
Gọi $r(u,v)$ là dung lượng trong mạng dư, khi đó
\[
  r(u,v)=c(u,v)-f(u,v).
\]
Nói cách khác, $r(u,v)$ là lượng lưu lượng còn có thể đẩy thêm qua cung
$(u,v)$. Nếu dung lượng của một cạnh trong mạng dư bằng $0$, ta xem như cạnh
đó không có trong mạng dư.

Ví dụ, nếu trên cạnh $s\to v_1$ có $(f,c)=(4,4)$, thì
$r(s,v_1)=0$ nên cạnh này không được vẽ trong mạng dư. Ngược lại, cạnh lùi
$v_1\to s$ có
\[
  r(v_1,s)=c(v_1,s)-f(v_1,s)=0-(-f(s,v_1))=4.
\]
Từ mạng dư ta đọc được ngay các khả năng tăng luồng: nếu có cạnh dư
$(s,v_2)$ dung lượng $3$, thì từ $s$ đến $v_2$ còn có thể tăng $3$ đơn vị;
nếu có cạnh dư $(v_1,t)$ dung lượng $2$, thì từ $v_1$ đến $t$ còn có thể tăng
$2$ đơn vị.

\subsection{Vì sao cần cung lùi?}

Cạnh như $(v_1,s)$ được gọi là cung lùi; nó biểu diễn rằng ta có thể tăng
$4$ đơn vị trên hướng $v_1\to s$ trong mạng dư. Thoạt nhìn, phát biểu ``từ
$v_1$ về $s$ còn có thể tăng $4$ đơn vị'' có vẻ vô nghĩa vì nó ngược hướng
cạnh trong mạng ban đầu. Tuy nhiên cung lùi chính là cơ chế cho phép thuật
toán sửa sai.

Trong ví dụ của slide, luồng đang vẽ chưa phải luồng cực đại. Nếu tăng thêm
$2$ đơn vị theo đường trong mạng dư
\[
  s\to v_2\to v_1\to t,
\]
ta thu được luồng cực đại. Đường này đi qua cung lùi $(v_2,v_1)$. Nếu không
tạo cung lùi, thuật toán sẽ không thể phát hiện đường tăng này. Về bản chất,
cung lùi cho phép thuật toán hủy một lựa chọn trước đó, ví dụ hủy việc đã cho
$2$ đơn vị luồng đi từ $v_1$ sang $v_2$.

\section{Đường tăng và thuật toán đường tăng}

\term{Đường cải thiện}, hay \term{đường tăng}, là một đường từ $s$ đến $t$ trong
mạng dư sao cho mọi cung $(u,v)$ trên đường đều có $r(u,v)>0$. Tương đương,
mọi cung tiến trên đường chưa bão hòa, và mọi cung lùi trên đường có luồng
khác không.

Thuật toán đường tăng làm như sau: mỗi lần dùng BFS tìm một đường tăng, rồi
điều chỉnh lưu lượng trên đường đó, thực tế là điều chỉnh dung lượng các cạnh
trong mạng dư, để tổng giá trị luồng tăng lên. Cách điều chỉnh là:

\begin{enumerate}
  \item Các cung không thuộc đường tăng $P$ giữ nguyên.
  \item Với mọi cung tiến trên $P$, cộng thêm $a$; với mọi cung lùi trên $P$,
  trừ đi $a$. Đại lượng $a$ là lượng tăng, được chọn lớn nhất có thể nhưng vẫn
  bảo đảm sau điều chỉnh
  \[
    0\le F_{ij}\le C_{ij}.
  \]
  Vì vậy
  \[
    a=\min\left(
      \min_{(i,j)\in P^+}(C_{ij}-F_{ij}),
      \min_{(i,j)\in P^-}F_{ij}
    \right),
  \]
  trong đó $P^+$ là tập cung tiến và $P^-$ là tập cung lùi.
\end{enumerate}

Nếu không còn đường tăng từ $s$ đến $t$, thuật toán dừng; khi đó luồng hiện
tại là luồng cực đại.

Quy trình tổng quát là: bắt đầu với $F_{ij}=0$ cho mọi cạnh; trong khi luồng
hiện tại còn đường tăng, thay nó bằng một luồng tốt hơn theo cách trên; nếu
không còn đường tăng thì kết luận $F$ là luồng cực đại. Câu hỏi còn lại là:
làm sao tìm đường tăng nhanh?

\section{Lát cắt và định lý luồng cực đại--lát cắt nhỏ nhất}

\subsection{Định nghĩa lát cắt}

Một lát cắt $(S,T)$ gồm hai tập đỉnh $S,T$ sao cho
\[
  S\cup T=V,\qquad s\in S,\qquad t\in T.
\]
Một cách xây dựng tập $S$ là bắt đầu từ $s\in S$ rồi mở rộng:
\begin{itemize}
  \item nếu $V_i\in S$ và $F_{ij}<C_{ij}$ thì đưa $V_j$ vào $S$;
  \item nếu $V_i\in S$ và $F_{ji}>0$ thì đưa $V_j$ vào $S$.
\end{itemize}
Một khi $t$ được đưa vào $S$, ta đã tìm thấy một đường tăng. Nếu $S$ không thể
mở rộng nữa mà $t$ vẫn chưa vào $S$, thì không tồn tại đường tăng; mọi đỉnh
ngoài $S$ được đưa vào $T$.

Với lát cắt $(S,T)$, lưu lượng qua lát cắt là
\[
  f(S,T)=\sum_{x\in S}\sum_{y\in T} f(x,y).
\]
Các hàm $c$ và $r$ cũng được định nghĩa tương tự, lần lượt là dung lượng của
lát cắt và dung lượng dư qua lát cắt.

\term{Định lý luồng cực đại--lát cắt nhỏ nhất.} Trong một mạng $D$, giá trị
luồng cực đại từ $s$ đến $t$ bằng dung lượng nhỏ nhất của một lát cắt tách
$s$ khỏi $t$.

\subsection{Ba điều kiện tương đương}

Trong một mạng luồng, ba điều kiện sau là tương đương tại cùng một thời điểm:

\begin{enumerate}
  \item $f$ là luồng cực đại.
  \item Trong mạng dư không còn đường tăng.
  \item Tồn tại lát cắt $(S,T)$ sao cho $|f|=c(S,T)$.
\end{enumerate}

Ta dùng một số kết luận phụ.

\begin{enumerate}
  \item $f(X,X)=0$, do phản đối xứng.
  \item $f(X,Y)=-f(Y,X)$, do phản đối xứng.
  \item $f(X\cup Y,Z)=f(X,Z)+f(Y,Z)$.
  \item $f(X,Y\cup Z)=f(X,Y)+f(X,Z)$.
\end{enumerate}

\term{Chứng minh $1\Rightarrow 2$.} Điều này hiển nhiên: nếu còn đường tăng,
vì dung lượng tăng ít nhất là dương, sau khi tăng theo đường đó ta thu được
một luồng có giá trị lớn hơn $f$, mâu thuẫn với việc $f$ là luồng cực đại.

\term{Kết luận phụ: tổng lưu lượng của một tập đỉnh bằng không.} Nếu $X$ không
chứa $s$ và $t$, thì tổng lưu lượng trên các cạnh liên thuộc với $X$ bằng
$0$:
\[
  f(X,V)=\sum_{x\in X}\left[\sum_{v\in V} f(x,v)\right]
        =\sum_{x\in X}0=0,
\]
do điều kiện bảo toàn luồng.

\term{Kết luận phụ: lưu lượng qua mọi lát cắt bằng giá trị luồng.} Với một lát
cắt $(S,T)$,
\[
\begin{aligned}
  f(S,T)
    &=f(S,V)-f(S,S)\\
    &=f(S,V)\\
    &=f(\{s\},V)+f(S\setminus\{s\},V)\\
    &=f(s,V)=|f|.
\end{aligned}
\]

\term{Kết luận phụ: giá trị luồng không vượt quá dung lượng của bất kỳ lát
cắt nào.} Cụ thể,
\[
\begin{aligned}
  |f|
    &=f(S,T)
      =\sum_{x\in S}\sum_{y\in T} f(x,y)\\
    &\le \sum_{x\in S}\sum_{y\in T} c(x,y)
      =c(S,T).
\end{aligned}
\]

\term{Chứng minh $2\Rightarrow 3$.} Đặt
\[
  S=\{s\}\cup\{v\mid \text{trong mạng dư có đường đi từ }s\text{ đến }v\},
  \qquad T=V\setminus S.
\]
Khi đó $(S,T)$ là một lát cắt. Theo kết luận phụ, $|f|=f(S,T)$. Hơn nữa
$r(S,T)=0$: nếu không, trong mạng dư sẽ có một cạnh từ một đỉnh $v\in S$ đến
một đỉnh $u\in T$, suy ra $s$ đi được đến $u$ qua $v$, trái với định nghĩa
của $S$. Do đó
\[
  |f|=f(S,T)=c(S,T)-r(S,T)=c(S,T).
\]

\term{Chứng minh $3\Rightarrow 1$.} Với mọi lát cắt ta luôn có
$|f|\le c(S,T)$. Nếu đã có $|f|=c(S,T)$ mà vẫn tồn tại luồng lớn hơn
$|f|+d$ với $d>0$, thì luồng đó sẽ vi phạm bất đẳng thức trên đối với cùng lát
cắt. Vì vậy $f$ là luồng cực đại.

\section{Phương pháp gán nhãn tìm đường tăng}

Phương pháp gán nhãn, tức thuật toán Ford--Fulkerson ở dạng cổ điển, bắt đầu
từ một luồng khả thi $F$; có thể lấy luồng không. Mỗi vòng gồm hai giai đoạn:
gán nhãn và điều chỉnh.

\subsection{Giai đoạn gán nhãn}

Mỗi đỉnh của mạng hoặc là đỉnh đã gán nhãn, hoặc chưa gán nhãn. Trong các đỉnh
đã gán nhãn lại chia thành đã kiểm tra và chưa kiểm tra. Nhãn của một đỉnh có
hai phần: đỉnh mà từ đó nhãn được truyền sang, và lượng tăng chắc chắn còn có
thể đưa tới đỉnh này.

Bắt đầu bằng cách gán cho nguồn $V_s$ nhãn $(0,+\infty)$; khi đó $V_s$ đã gán
nhãn nhưng chưa kiểm tra, còn mọi đỉnh khác chưa có nhãn. Lấy một đỉnh
$V_i$ đã gán nhãn nhưng chưa kiểm tra, rồi xét mọi đỉnh $V_j$ chưa gán nhãn:

\begin{enumerate}
  \item Nếu có cung tiến $(V_i,V_j)$ và $F_{ij}<C_{ij}$, gán cho $V_j$ nhãn
  \[
    (V_i,L(V_j)),\qquad
    L(V_j)=\min\{L(V_i),\,C_{ij}-F_{ij}\}.
  \]
  Khi đó $V_j$ trở thành đỉnh đã gán nhãn nhưng chưa kiểm tra.

  \item Nếu có cung lùi tương ứng với $(V_j,V_i)$ và $F_{ji}>0$, gán cho
  $V_j$ nhãn
  \[
    (-V_i,L(V_j)),\qquad
    L(V_j)=\min\{L(V_i),\,F_{ji}\}.
  \]
  Dấu trừ cho biết bước này đi qua một cung lùi.
\end{enumerate}

Sau khi mọi đỉnh kề $V_i$ đều đã được xử lý, $V_i$ trở thành đỉnh đã kiểm tra.
Lặp lại quá trình trên. Một khi $V_t$ được gán nhãn, ta đã tìm được một đường
tăng từ $V_s$ đến $V_t$ và chuyển sang giai đoạn điều chỉnh.

\subsection{Giai đoạn điều chỉnh}

Ta truy ngược từ $V_t$ theo phần đầu của nhãn để lấy đường tăng $P$. Nếu nhãn
đầu của $V_t$ là $V_k$ hoặc $-V_k$, thì cung tương ứng trên $P$ là
$(V_k,V_t)$ hoặc cung lùi của $(V_t,V_k)$. Sau đó tiếp tục xem nhãn của
$V_k$, cứ thế cho đến khi truy ngược về $V_s$.

Lượng tăng là
\[
  a=L(V_t).
\]
Luồng được cập nhật theo công thức
\[
  F'_{ij}=
  \begin{cases}
    F_{ij}+a, & (V_i,V_j)\in P^+,\\
    F_{ij}-a, & (V_i,V_j)\in P^-,\\
    F_{ij},   & (V_i,V_j)\notin P.
  \end{cases}
\]
Xóa toàn bộ nhãn rồi lặp lại giai đoạn gán nhãn với luồng mới $F'$. Thuật toán
dừng khi gán nhãn không thể tiếp tục.

\subsection{Ví dụ gán nhãn}

Ví dụ trong slide có các đỉnh $V_s,V_1,V_2,V_3,V_4,V_t$. Mỗi cạnh dưới đây
được ghi dưới dạng $(f,c)$:
\[
\begin{array}{c|c}
\text{Cạnh} & (f,c)\\ \hline
V_s\to V_2 & (3,3)\\
V_s\to V_1 & (1,5)\\
V_2\to V_4 & (3,4)\\
V_2\to V_1 & (1,1)\\
V_2\to V_3 & (1,1)\\
V_1\to V_3 & (2,2)\\
V_4\to V_t & (3,5)\\
V_3\to V_t & (1,2)\\
V_3\to V_4 & (0,3)
\end{array}
\]

Quá trình gán nhãn:
\begin{enumerate}
  \item Đầu tiên gán cho $V_s$ nhãn $(0,+\infty)$.
  \item Kiểm tra $V_s$. Trên cung $V_s\to V_2$ có $F_{s2}=C_{s2}=3$, nên
  không thể gán nhãn qua cung này. Trên cung $V_s\to V_1$ có
  $F_{s1}=1<C_{s1}=5$, nên gán cho $V_1$ nhãn
  \[
    (V_s,L(V_1)),\qquad L(V_1)=\min\{+\infty,5-1\}=4.
  \]
  \item Kiểm tra $V_1$. Cung $V_1\to V_3$ đã bão hòa vì $F_{13}=C_{13}=2$.
  Nhưng trên cung $V_2\to V_1$ có $F_{21}=1>0$, nên có thể đi lùi từ $V_1$
  tới $V_2$ và gán cho $V_2$ nhãn
  \[
    (-V_1,L(V_2)),\qquad L(V_2)=\min\{4,1\}=1.
  \]
  \item Kiểm tra $V_2$. Vì $F_{24}=3<C_{24}=4$, gán cho $V_4$ nhãn
  \[
    (V_2,L(V_4)),\qquad L(V_4)=\min\{1,1\}=1.
  \]
  Tương tự, qua cung lùi liên quan tới $V_3$ ta gán được cho $V_3$ nhãn
  $(-V_2,1)$.
  \item Chọn một trong hai đỉnh $V_3,V_4$ để kiểm tra, chẳng hạn $V_4$. Trên
  cung $V_4\to V_t$ có $F_{4t}=3<C_{4t}=5$, nên gán cho $V_t$ nhãn
  \[
    (V_4,L(V_t)),\qquad L(V_t)=\min\{1,2\}=1.
  \]
\end{enumerate}

Vì $V_t$ đã có nhãn, ta điều chỉnh. Đường tăng tìm được bằng truy ngược nhãn
có các cung tiến
\[
  P^+=\{(V_s,V_1),(V_3,V_t)\}
\]
và các cung lùi
\[
  P^-=\{(V_2,V_1),(V_3,V_2)\}.
\]
Với $a=1$, cập nhật
\[
  F'_{s1}=F_{s1}+1=2,\qquad F'_{3t}=F_{3t}+1=2,
\]
\[
  F'_{21}=F_{21}-1=0,\qquad F'_{32}=F_{32}-1=0.
\]

Sau khi gán nhãn lại, khi kiểm tra đến $V_1$ thì $F_{21}=0$ và
$F_{13}=C_{13}$ đều không thỏa điều kiện gán nhãn, nên quá trình không thể
tiếp tục. Luồng khả thi khi đó có giá trị cực đại bằng $5$. Đồng thời ta tìm
được lát cắt nhỏ nhất
\[
  S=\{V_s,V_1\},\qquad T=\{V_2,V_3,V_4,V_t\},
\]
với
\[
  C(S,T)=5.
\]
Điều này kiểm chứng định lý luồng cực đại--lát cắt nhỏ nhất.

\section{Mạng nhiều nguồn, nhiều đích}

Mạng nhiều nguồn, nhiều đích có thể chuyển về bài toán một nguồn, một đích.
Thêm một siêu nguồn nối tới mọi nguồn ban đầu bằng cạnh dung lượng đủ lớn, và
thêm một siêu đích nhận cạnh từ mọi đích ban đầu, cũng với dung lượng đủ lớn.
Sau đó giải bài toán luồng cực đại thông thường giữa siêu nguồn và siêu đích.

\section{Luồng cực đại với cận dưới và cận trên}

Trong mạng có cận dưới và cận trên, mỗi cung $e$ đi kèm hai số $B(e)$ và
$C(e)$, lần lượt là cận dưới và cận trên của lưu lượng. Ta cần tìm luồng cực
đại thỏa
\[
  B(e)\le F(e)\le C(e).
\]
Mạng trong phương pháp gán nhãn ở trên là trường hợp đặc biệt với $B(e)=0$.
Khi có cạnh với $B(e)>0$, mạng chưa chắc có luồng khả thi. Vì vậy trước hết
phải biết cách kiểm tra sự tồn tại của luồng khả thi.

\subsection{Chuyển thành mạng phụ}

Ý tưởng là biến mạng ban đầu thành một mạng phụ.

\begin{enumerate}
  \item Thêm hai đỉnh mới $s^*$ và $t^*$, gọi là nguồn phụ và đích phụ.
  \item Với mỗi đỉnh $U$ của mạng ban đầu, thêm cung $U\to t^*$ có dung lượng
  bằng tổng các cận dưới của những cung đi ra khỏi $U$.
  \item Với mỗi đỉnh $U$ của mạng ban đầu, thêm cung $s^*\to U$ có dung lượng
  bằng tổng các cận dưới của những cung đi vào $U$.
  \item Giữ các cung của mạng ban đầu, nhưng đổi dung lượng của cung $e$ thành
  $C(e)-B(e)$.
  \item Thêm hai cung giữa nguồn và đích ban đầu theo hai chiều, dung lượng
  đều là $+\infty$.
  \item Dùng phương pháp gán nhãn tìm luồng cực đại trong mạng mới. Nếu mọi
  cung đi ra khỏi $s^*$ đều bão hòa, thì mạng ban đầu có luồng khả thi; khi đó
  khôi phục bằng
  \[
    F(e)=F'(e)+B(e).
  \]
  Nếu không, mạng ban đầu vô nghiệm.
  \item Trên mạng ban đầu, tiếp tục dùng phương pháp gán nhãn để tăng luồng
  khả thi vừa tìm được, từ đó thu được luồng cực đại.
\end{enumerate}

Trong ví dụ của slide, mạng ban đầu ghi mỗi cạnh bằng cặp $(B(e),C(e))$:
\[
\begin{array}{c|c}
\text{Cạnh} & (B(e),C(e))\\ \hline
s\to w & (1,3)\\
s\to y & (0,10)\\
w\to x & (5,7)\\
x\to t & (3,5)\\
y\to z & (2,8)\\
z\to t & (2,6)\\
x\to y & (1,3)\\
z\to w & (2,4)
\end{array}
\]
Theo thuật toán trên, ta dựng mạng phụ và chạy gán nhãn. Kết quả cho thấy mọi
cung đi ra khỏi nguồn phụ đều bão hòa, nên tồn tại luồng khả thi. Sau khi khôi
phục theo công thức $F(e)=F'(e)+B(e)$ và tiếp tục tăng luồng khả thi trên mạng
ban đầu, ta thu được luồng cực đại có giá trị $10$.

\subsection{Luồng nhỏ nhất với cận dưới và cận trên}

Để tìm luồng nhỏ nhất, trước hết cũng dùng phương pháp mạng phụ để tìm một
luồng khả thi. Khác biệt nằm ở bước tăng luồng: thay vì tiếp tục tăng từ
nguồn ban đầu $s$ đến đích ban đầu $t$, ta lấy $t$ làm nguồn và $s$ làm đích
để tìm luồng cực đại theo chiều ngược. Lượng luồng đẩy ngược được trừ khỏi
luồng khả thi ban đầu; kết quả là luồng nhỏ nhất từ $s$ đến $t$.

\section{Luồng cực đại chi phí nhỏ nhất}

Bài toán luồng cực đại chi phí nhỏ nhất yêu cầu tìm phương án vận chuyển có
lượng vận chuyển lớn nhất, và trong các phương án đó có tổng chi phí nhỏ nhất.
Nếu $b_{ij}$ là chi phí đơn vị trên cung $(i,j)$, tổng chi phí của luồng là
\[
  B(F)=\sum b_{ij}F_{ij}.
\]
Ta cần tìm một luồng cực đại làm biểu thức này nhỏ nhất.

\subsection{Tư tưởng thuật toán}

Nếu $F$ là luồng có chi phí nhỏ nhất trong tất cả các luồng khả thi có cùng
giá trị hiện tại, và $P$ là đường tăng có chi phí nhỏ nhất đối với $F$, thì
sau khi tăng luồng tối đa dọc theo $P$, luồng mới vẫn là luồng có chi phí nhỏ
nhất trong lớp giá trị mới. Lặp lại quá trình này cho tới khi không còn đường
tăng sẽ thu được luồng cực đại chi phí nhỏ nhất.

\subsection{Các bước}

\begin{enumerate}
  \item Lấy $F^{(0)}=0$.
  \item Giả sử sau bước $k-1$ ta đã có luồng chi phí nhỏ nhất $F^{(k-1)}$.
  Dựng đồ thị có trọng số $W(F^{(k-1)})$, rồi tìm đường ngắn nhất từ $V_s$ đến
  $V_t$ trong đồ thị này. Nếu không tồn tại đường ngắn nhất, tức khoảng cách
  là $+\infty$, thì $F^{(k-1)}$ là luồng cực đại chi phí nhỏ nhất. Nếu tồn tại,
  đường ngắn nhất đó là đường tăng $P$, và ta điều chỉnh $F^{(k-1)}$ trên $P$.
  \item Cách dựng $W(F)$: tập đỉnh giữ nguyên như mạng ban đầu. Mỗi cung
  $(V_i,V_j)$ của mạng ban đầu được xem xét theo hai hướng. Trọng số hướng
  thuận là
  \[
    W_{ij}=
    \begin{cases}
      b_{ij}, & F_{ij}<C_{ij},\\
      +\infty, & F_{ij}=C_{ij},
    \end{cases}
  \]
  còn trọng số hướng ngược là
  \[
    W_{ji}=
    \begin{cases}
      -b_{ij}, & F_{ij}>0,\\
      +\infty, & F_{ij}=0.
    \end{cases}
  \]
  Các cạnh có độ dài $+\infty$ có thể bỏ khỏi $W(F)$.
  \item Nếu $P^+$ là tập cung thuận và $P^-$ là tập cung lùi trên đường tăng,
  lượng tăng là
  \[
    a=\min\left[
      \min_{(i,j)\in P^+}\bigl(C_{ij}-F_{ij}^{(k-1)}\bigr),
      \min_{(i,j)\in P^-}F_{ij}^{(k-1)}
    \right].
  \]
  Cập nhật
  \[
    F_{ij}^{(k)}=
    \begin{cases}
      F_{ij}^{(k-1)}+a, & (V_i,V_j)\in P^+,\\
      F_{ij}^{(k-1)}-a, & (V_i,V_j)\in P^-,\\
      F_{ij}^{(k-1)}, & (V_i,V_j)\notin P.
    \end{cases}
  \]
  \item Với luồng mới $F^{(k)}$, lặp lại các bước trên cho tới khi không còn
  đường ngắn nhất từ nguồn đến đích.
\end{enumerate}

\subsection{Ví dụ}

Ví dụ trong slide dùng cặp $(C,b)$, trong đó $C$ là dung lượng và $b$ là chi
phí đơn vị:
\[
\begin{array}{c|c}
\text{Cạnh} & (C,b)\\ \hline
V_s\to V_1 & (10,4)\\
V_s\to V_2 & (8,1)\\
V_2\to V_1 & (5,2)\\
V_1\to V_t & (7,1)\\
V_2\to V_3 & (10,3)\\
V_1\to V_3 & (2,6)\\
V_3\to V_t & (4,2)
\end{array}
\]

Với $F^{(0)}=0$, đồ thị $W(F^{(0)})$ chỉ có các cạnh thuận có chi phí
$4,1,2,1,3,6,2$ tương ứng với bảng trên. Đường ngắn nhất đầu tiên là
\[
  V_s\to V_2\to V_1\to V_t,
\]
có lượng tăng $5$. Ta thu được $F^{(1)}$ với giá trị luồng $v(F^{(1)})=5$:
\[
\begin{array}{c|c}
\text{Cạnh} & F^{(1)}\\ \hline
V_s\to V_1 & 0\\
V_s\to V_2 & 5\\
V_2\to V_1 & 5\\
V_1\to V_t & 5\\
V_2\to V_3 & 0\\
V_1\to V_3 & 0\\
V_3\to V_t & 0
\end{array}
\]

Trong $W(F^{(1)})$, các cạnh đã có luồng sinh ra cạnh ngược chi phí âm. Đường
ngắn nhất kế tiếp là $V_s\to V_1\to V_t$, tăng được $2$ đơn vị, nên
$v(F^{(2)})=7$:
\[
\begin{array}{c|c}
\text{Cạnh} & F^{(2)}\\ \hline
V_s\to V_1 & 2\\
V_s\to V_2 & 5\\
V_2\to V_1 & 5\\
V_1\to V_t & 7\\
V_2\to V_3 & 0\\
V_1\to V_3 & 0\\
V_3\to V_t & 0
\end{array}
\]

Lần tiếp theo, đường ngắn nhất là
\[
  V_s\to V_2\to V_3\to V_t,
\]
tăng được $3$ đơn vị và cho $v(F^{(3)})=10$:
\[
\begin{array}{c|c}
\text{Cạnh} & F^{(3)}\\ \hline
V_s\to V_1 & 2\\
V_s\to V_2 & 8\\
V_2\to V_1 & 5\\
V_1\to V_t & 7\\
V_2\to V_3 & 3\\
V_1\to V_3 & 0\\
V_3\to V_t & 3
\end{array}
\]

Cuối cùng, thuật toán dùng một đường có cung lùi:
\[
  V_s\to V_1\to V_2\to V_3\to V_t,
\]
trong đó bước $V_1\to V_2$ là cạnh ngược của cung $V_2\to V_1$. Lượng tăng là
$1$, thu được $v(F^{(4)})=11$:
\[
\begin{array}{c|c}
\text{Cạnh} & F^{(4)}\\ \hline
V_s\to V_1 & 3\\
V_s\to V_2 & 8\\
V_2\to V_1 & 4\\
V_1\to V_t & 7\\
V_2\to V_3 & 4\\
V_1\to V_3 & 0\\
V_3\to V_t & 4
\end{array}
\]
Vì tổng dung lượng vào đích là $7+4=11$, đây đã là luồng cực đại. Do mỗi bước
đều chọn đường tăng ngắn nhất theo chi phí trong mạng dư, luồng cực đại này
cũng có chi phí nhỏ nhất.

\section{Bài tập}

\begin{itemize}
  \item Cài đặt các thuật toán trong bài.
  \item Phương pháp gán nhãn.
  \item Luồng cực đại trong mạng có cận dưới và cận trên.
  \item Luồng nhỏ nhất trong mạng có cận dưới và cận trên.
  \item Luồng cực đại chi phí nhỏ nhất.
\end{itemize}

\end{document}
